\section{분배 전략과 대화 분석}\label{sec:strategies}

\noindent\textcolor{red}{\textbf{예상!}} 분석 후 작성할 내용을 위해 남겨 둔 공간이다.

\subsection{분배 시점과 금액}

매 라운드 잔량을 맞추는 분배와 오답 뒤 필요한 만큼 보충하는 분배를 비교한다.
잦은 소액 이전과 적은 횟수의 큰 이전이 어떻게 다른지도 살펴본다.

\subsection{대화에 나타난 판단}

응답과 대화에서 말한 분배 이유를 실제 토큰 이동과 대조한다.

\begin{figure*}[t]
\centering
\resizebox{\textwidth}{!}{%% Generated by cogsci2027/tools/arena_gini.py from results/arena_v33_gini.json. Do not edit by hand.
%% v33, 5 games per model (seeds s4101, s4102, s4103, s4104, s4105); largest single-round total given 37,600 tokens.
\begin{tikzpicture}
\begin{groupplot}[sgaxis,
  group style={group size=4 by 3,horizontal sep=0.35cm,vertical sep=0.28cm,
    y descriptions at=edge left,x descriptions at=edge bottom},
  width=4.75cm,xmin=.9,xmax=9.1,xtick={1.5,2.5,3.5,4.5,5.5,6.5,7.5,8.5},
  xticklabels={},xtick style={draw=none},tick label style={font=\tiny},
  xticklabel style={align=center},title style={font=\small,align=center}]
\nextgroupplot[height=3.8cm,ymin=0,ymax=.75,ytick={0,.25,.5,.75},ylabel={지니 계수},title={GPT-6 Astra\\[-1pt]\scriptsize Sen 3.52M $\cdot$ 끝까지 4/5판}]
\addplot[sgastra,opacity=.22,line width=.4pt,forget plot,unbounded coords=jump] coordinates {(1.00,0.3269) (1.25,0.3269) (1.50,0.3373) (1.75,0.2907) (2.00,0.3590) (2.25,0.3590) (2.50,0.3651) (2.75,0.2586) (3.00,0.3390) (3.25,0.3390) (3.50,0.2487) (3.75,0.2487) (4.00,0.2488) (4.25,0.2488) (4.50,0.2490) (4.75,0.2490) (5.00,0.2490) (5.25,0.2490) (5.50,0.2528) (5.75,0.2528) (6.00,0.3826) (6.25,0.3826) (nan,nan) (nan,nan) (nan,nan) (nan,nan) (nan,nan) (nan,nan) (nan,nan) (nan,nan) (nan,nan) (nan,nan)};
\addplot[sgastra,opacity=.22,line width=.4pt,forget plot,unbounded coords=jump] coordinates {(1.00,0.3269) (1.25,0.3269) (1.50,0.3277) (1.75,0.3277) (2.00,0.3278) (2.25,0.3278) (2.50,0.3278) (2.75,0.3278) (3.00,0.3279) (3.25,0.3279) (3.50,0.3355) (3.75,0.2423) (4.00,0.2986) (4.25,0.2986) (4.50,0.2987) (4.75,0.2987) (5.00,0.2987) (5.25,0.2987) (5.50,0.2988) (5.75,0.2988) (6.00,0.2989) (6.25,0.2989) (6.50,0.2957) (6.75,0.1367) (7.00,0.1777) (7.25,0.1777) (7.50,0.1778) (7.75,0.1778) (8.00,0.1778) (8.25,0.1778) (8.50,0.1779) (8.75,0.1779)};
\addplot[sgastra,opacity=.22,line width=.4pt,forget plot,unbounded coords=jump] coordinates {(1.00,0.3269) (1.25,0.3269) (1.50,0.3271) (1.75,0.3271) (2.00,0.3271) (2.25,0.3271) (2.50,0.3272) (2.75,0.3272) (3.00,0.3272) (3.25,0.3272) (3.50,0.3273) (3.75,0.3273) (4.00,0.3273) (4.25,0.3273) (4.50,0.3274) (4.75,0.3274) (5.00,0.3274) (5.25,0.3274) (5.50,0.3247) (5.75,0.3247) (6.00,0.3989) (6.25,0.3989) (6.50,0.2585) (6.75,0.2585) (7.00,0.2585) (7.25,0.2585) (7.50,0.2609) (7.75,0.2609) (8.00,0.3435) (8.25,0.3435) (8.50,0.1397) (8.75,0.1397)};
\addplot[sgastra,opacity=.22,line width=.4pt,forget plot,unbounded coords=jump] coordinates {(1.00,0.3269) (1.25,0.3269) (1.50,0.3274) (1.75,0.3274) (2.00,0.3275) (2.25,0.2119) (2.50,0.2172) (2.75,0.2172) (3.00,0.2675) (3.25,0.2675) (3.50,0.2676) (3.75,0.2676) (4.00,0.2676) (4.25,0.2676) (4.50,0.2677) (4.75,0.2677) (5.00,0.2677) (5.25,0.2677) (5.50,0.2678) (5.75,0.2678) (6.00,0.2678) (6.25,0.2678) (6.50,0.2679) (6.75,0.2679) (7.00,0.2679) (7.25,0.2679) (7.50,0.2680) (7.75,0.2680) (8.00,0.2680) (8.25,0.2680) (8.50,0.2681) (8.75,0.2681)};
\addplot[sgastra,opacity=.22,line width=.4pt,forget plot,unbounded coords=jump] coordinates {(1.00,0.3269) (1.25,0.3269) (1.50,0.3275) (1.75,0.3275) (2.00,0.3276) (2.25,0.3276) (2.50,0.3277) (2.75,0.3277) (3.00,0.3277) (3.25,0.3277) (3.50,0.3310) (3.75,0.2379) (4.00,0.2930) (4.25,0.2930) (4.50,0.2930) (4.75,0.2930) (5.00,0.2931) (5.25,0.2931) (5.50,0.2931) (5.75,0.2931) (6.00,0.2932) (6.25,0.2932) (6.50,0.2932) (6.75,0.2932) (7.00,0.2933) (7.25,0.2933) (7.50,0.2937) (7.75,0.1804) (8.00,0.2355) (8.25,0.2355) (8.50,0.2355) (8.75,0.2355)};
\addplot[black!45,densely dotted,line width=.7pt,forget plot,unbounded coords=jump] coordinates {(1.00,0.3269) (1.25,0.3269) (1.50,0.3294) (1.75,0.3201) (2.00,0.3338) (2.25,0.3107) (2.50,0.3130) (2.75,0.2917) (3.00,0.3178) (3.25,0.3178) (3.50,0.3020) (3.75,0.2648) (4.00,0.2871) (4.25,0.2871) (4.50,0.2871) (4.75,0.2871) (5.00,0.2872) (5.25,0.2872) (5.50,0.2874) (5.75,0.2874) (6.00,0.3283) (6.25,0.3283) (6.50,0.2788) (6.75,0.2391) (7.00,0.2494) (7.25,0.2494) (7.50,0.2501) (7.75,0.2218) (8.00,0.2562) (8.25,0.2562) (8.50,0.2053) (8.75,0.2053)};
\addplot[sgastra,line width=1.4pt,forget plot,unbounded coords=jump] coordinates {(1.00,0.3269) (1.25,0.3269) (nan,nan) (1.50,0.3294) (1.75,0.3201) (nan,nan) (2.00,0.3338) (2.25,0.3107) (nan,nan) (2.50,0.3130) (2.75,0.2917) (nan,nan) (3.00,0.3178) (3.25,0.3178) (nan,nan) (3.50,0.3020) (3.75,0.2648) (nan,nan) (4.00,0.2871) (4.25,0.2871) (nan,nan) (4.50,0.2871) (4.75,0.2871) (nan,nan) (5.00,0.2872) (5.25,0.2872) (nan,nan) (5.50,0.2874) (5.75,0.2874) (nan,nan) (6.00,0.3283) (6.25,0.3283) (nan,nan) (6.50,0.2788) (6.75,0.2391) (nan,nan) (7.00,0.2494) (7.25,0.2494) (nan,nan) (7.50,0.2501) (7.75,0.2218) (nan,nan) (8.00,0.2562) (8.25,0.2562) (nan,nan) (8.50,0.2053) (8.75,0.2053) (nan,nan)};
\nextgroupplot[height=3.8cm,ymin=0,ymax=.75,ytick={0,.25,.5,.75},title={GPT-6 Sol\\[-1pt]\scriptsize Sen 3.02M $\cdot$ 끝까지 4/5판}]
\addplot[sgsol,opacity=.22,line width=.4pt,forget plot,unbounded coords=jump] coordinates {(1.00,0.3269) (1.25,0.3269) (1.50,0.3432) (1.75,0.3432) (2.00,0.4416) (2.25,0.4416) (2.50,0.3014) (2.75,0.3014) (3.00,0.3924) (3.25,0.3924) (3.50,0.1408) (3.75,0.1408) (4.00,0.1420) (4.25,0.1420) (4.50,0.1453) (4.75,0.1453) (5.00,0.1498) (5.25,0.1498) (5.50,0.1716) (5.75,0.1716) (6.00,0.2842) (6.25,0.2842) (nan,nan) (nan,nan) (nan,nan) (nan,nan) (nan,nan) (nan,nan) (nan,nan) (nan,nan) (nan,nan) (nan,nan)};
\addplot[sgsol,opacity=.22,line width=.4pt,forget plot,unbounded coords=jump] coordinates {(1.00,0.3269) (1.25,0.3269) (1.50,0.3325) (1.75,0.3325) (2.00,0.3350) (2.25,0.3350) (2.50,0.3391) (2.75,0.3391) (3.00,0.3409) (3.25,0.3409) (3.50,0.3623) (3.75,0.3341) (4.00,0.4351) (4.25,0.3897) (4.50,0.4202) (4.75,0.4202) (5.00,0.4368) (5.25,0.4368) (5.50,0.4502) (5.75,0.4502) (6.00,0.4561) (6.25,0.4561) (6.50,0.3571) (6.75,0.3320) (7.00,0.4460) (7.25,0.4460) (7.50,0.4616) (7.75,0.3171) (8.00,0.4651) (8.25,0.4651) (8.50,0.4802) (8.75,0.4802)};
\addplot[sgsol,opacity=.22,line width=.4pt,forget plot,unbounded coords=jump] coordinates {(1.00,0.3269) (1.25,0.3269) (1.50,0.3278) (1.75,0.3278) (2.00,0.3296) (2.25,0.3296) (2.50,0.3313) (2.75,0.3313) (3.00,0.3330) (3.25,0.3330) (3.50,0.3390) (3.75,0.3390) (4.00,0.3401) (4.25,0.3401) (4.50,0.3442) (4.75,0.3442) (5.00,0.3465) (5.25,0.3465) (5.50,0.3461) (5.75,0.3461) (6.00,0.4507) (6.25,0.4507) (6.50,0.2786) (6.75,0.2786) (7.00,0.2795) (7.25,0.2795) (7.50,0.2859) (7.75,0.2859) (8.00,0.4042) (8.25,0.4042) (8.50,0.1573) (8.75,0.1573)};
\addplot[sgsol,opacity=.22,line width=.4pt,forget plot,unbounded coords=jump] coordinates {(1.00,0.3269) (1.25,0.3269) (1.50,0.3302) (1.75,0.3302) (2.00,0.3324) (2.25,0.3324) (2.50,0.3614) (2.75,0.2104) (3.00,0.2736) (3.25,0.2736) (3.50,0.2761) (3.75,0.2761) (4.00,0.2772) (4.25,0.2772) (4.50,0.2794) (4.75,0.2794) (5.00,0.2803) (5.25,0.2803) (5.50,0.2808) (5.75,0.2808) (6.00,0.2815) (6.25,0.2815) (6.50,0.2843) (6.75,0.2843) (7.00,0.2850) (7.25,0.2850) (7.50,0.2864) (7.75,0.2864) (8.00,0.2871) (8.25,0.2871) (8.50,0.2894) (8.75,0.2894)};
\addplot[sgsol,opacity=.22,line width=.4pt,forget plot,unbounded coords=jump] coordinates {(1.00,0.3269) (1.25,0.3269) (1.50,0.3321) (1.75,0.3321) (2.00,0.3364) (2.25,0.3364) (2.50,0.3376) (2.75,0.3376) (3.00,0.3398) (3.25,0.3398) (3.50,0.3426) (3.75,0.3426) (4.00,0.3452) (4.25,0.3452) (4.50,0.3505) (4.75,0.3505) (5.00,0.3529) (5.25,0.3529) (5.50,0.3560) (5.75,0.2829) (6.00,0.2868) (6.25,0.2868) (6.50,0.2873) (6.75,0.2873) (7.00,0.2878) (7.25,0.2878) (7.50,0.2903) (7.75,0.2648) (8.00,0.3412) (8.25,0.3412) (8.50,0.3438) (8.75,0.3438)};
\addplot[black!45,densely dotted,line width=.7pt,forget plot,unbounded coords=jump] coordinates {(1.00,0.3269) (1.25,0.3269) (1.50,0.3332) (1.75,0.3332) (2.00,0.3550) (2.25,0.3550) (2.50,0.3341) (2.75,0.3040) (3.00,0.3359) (3.25,0.3359) (3.50,0.2921) (3.75,0.2865) (4.00,0.3079) (4.25,0.2988) (4.50,0.3079) (4.75,0.3079) (5.00,0.3133) (5.25,0.3133) (5.50,0.3209) (5.75,0.3063) (6.00,0.3519) (6.25,0.3519) (6.50,0.3018) (6.75,0.2956) (7.00,0.3246) (7.25,0.3246) (7.50,0.3310) (7.75,0.2886) (8.00,0.3744) (8.25,0.3744) (8.50,0.3177) (8.75,0.3177)};
\addplot[sgsol,line width=1.4pt,forget plot,unbounded coords=jump] coordinates {(1.00,0.3269) (1.25,0.3269) (nan,nan) (1.50,0.3332) (1.75,0.3332) (nan,nan) (2.00,0.3550) (2.25,0.3550) (nan,nan) (2.50,0.3341) (2.75,0.3040) (nan,nan) (3.00,0.3359) (3.25,0.3359) (nan,nan) (3.50,0.2921) (3.75,0.2865) (nan,nan) (4.00,0.3079) (4.25,0.2988) (nan,nan) (4.50,0.3079) (4.75,0.3079) (nan,nan) (5.00,0.3133) (5.25,0.3133) (nan,nan) (5.50,0.3209) (5.75,0.3063) (nan,nan) (6.00,0.3519) (6.25,0.3519) (nan,nan) (6.50,0.3018) (6.75,0.2956) (nan,nan) (7.00,0.3246) (7.25,0.3246) (nan,nan) (7.50,0.3310) (7.75,0.2886) (nan,nan) (8.00,0.3744) (8.25,0.3744) (nan,nan) (8.50,0.3177) (8.75,0.3177) (nan,nan)};
\nextgroupplot[height=3.8cm,ymin=0,ymax=.75,ytick={0,.25,.5,.75},title={Opus 5.5\\[-1pt]\scriptsize Sen 2.27M $\cdot$ 끝까지 3/5판}]
\addplot[sgopus,opacity=.22,line width=.4pt,forget plot,unbounded coords=jump] coordinates {(1.00,0.3269) (1.25,0.2538) (1.50,0.2688) (1.75,0.2688) (2.00,0.3379) (2.25,0.3379) (2.50,0.3559) (2.75,0.2868) (3.00,0.3973) (3.25,0.3973) (3.50,0.1561) (3.75,0.1561) (4.00,0.1571) (4.25,0.1571) (4.50,0.1633) (4.75,0.1633) (5.00,0.1635) (5.25,0.1635) (5.50,0.1724) (5.75,0.1724) (6.00,0.3166) (6.25,0.3166) (nan,nan) (nan,nan) (nan,nan) (nan,nan) (nan,nan) (nan,nan)};
\addplot[sgopus,opacity=.22,line width=.4pt,forget plot,unbounded coords=jump] coordinates {(1.00,0.3269) (1.25,0.3269) (1.50,0.2252) (1.75,0.2252) (2.00,0.2258) (2.25,0.2258) (2.50,0.2267) (2.75,0.2267) (3.00,0.2273) (3.25,0.2273) (3.50,0.2259) (3.75,0.2259) (4.00,0.2861) (4.25,0.2861) (4.50,0.2963) (4.75,0.2963) (5.00,0.2972) (5.25,0.2972) (5.50,0.3044) (5.75,0.3044) (6.00,0.3047) (6.25,0.3047) (6.50,0.3121) (6.75,0.2206) (7.00,0.3232) (7.25,0.3232) (7.50,0.3350) (7.75,0.3350) (8.00,0.3405) (8.25,0.3405) (8.50,0.3574) (8.75,0.3972)};
\addplot[sgopus,opacity=.22,line width=.4pt,forget plot,unbounded coords=jump] coordinates {(1.00,0.3269) (1.25,0.2635) (1.50,0.2713) (1.75,0.2772) (2.00,0.2852) (2.25,0.2852) (2.50,0.2874) (2.75,0.2874) (3.00,0.2878) (3.25,0.2878) (3.50,0.2853) (3.75,0.2609) (4.00,0.3269) (4.25,0.3269) (4.50,0.3343) (4.75,0.2509) (5.00,0.3599) (5.25,0.2882) (5.50,0.3403) (5.75,0.2267) (6.00,0.3227) (6.25,0.2916) (6.50,0.1273) (6.75,0.1273) (7.00,0.1418) (7.25,0.2997) (7.50,0.4892) (7.75,0.4559)};
\addplot[sgopus,opacity=.22,line width=.4pt,forget plot,unbounded coords=jump] coordinates {(1.00,0.3269) (1.25,0.3269) (1.50,0.3332) (1.75,0.3332) (2.00,0.3335) (2.25,0.1985) (2.50,0.2071) (2.75,0.2609) (3.00,0.3391) (3.25,0.3391) (3.50,0.3512) (3.75,0.3512) (4.00,0.3526) (4.25,0.3526) (4.50,0.3605) (4.75,0.2521) (5.00,0.2520) (5.25,0.2520) (5.50,0.2541) (5.75,0.2541) (6.00,0.2565) (6.25,0.2565) (6.50,0.2601) (6.75,0.2601) (7.00,0.2600) (7.25,0.2600) (7.50,0.2632) (7.75,0.2632) (8.00,0.2653) (8.25,0.2653) (8.50,0.2684) (8.75,0.1765)};
\addplot[sgopus,opacity=.22,line width=.4pt,forget plot,unbounded coords=jump] coordinates {(1.00,0.3269) (1.25,0.3269) (1.50,0.2284) (1.75,0.2284) (2.00,0.2290) (2.25,0.2290) (2.50,0.2290) (2.75,0.2290) (3.00,0.2291) (3.25,0.2291) (3.50,0.2318) (3.75,0.2318) (4.00,0.2320) (4.25,0.2320) (4.50,0.2344) (4.75,0.2344) (5.00,0.2347) (5.25,0.2347) (5.50,0.2376) (5.75,0.2376) (6.00,0.2378) (6.25,0.2378) (6.50,0.2379) (6.75,0.2379) (7.00,0.2382) (7.25,0.2382) (7.50,0.2414) (7.75,0.2414) (8.00,0.2416) (8.25,0.2416) (8.50,0.2440) (8.75,0.2440)};
\addplot[black!45,densely dotted,line width=.7pt,forget plot,unbounded coords=jump] coordinates {(1.00,0.3269) (1.25,0.2996) (1.50,0.2654) (1.75,0.2666) (2.00,0.2823) (2.25,0.2553) (2.50,0.2612) (2.75,0.2582) (3.00,0.2961) (3.25,0.2961) (3.50,0.2501) (3.75,0.2452) (4.00,0.2709) (4.25,0.2709) (4.50,0.2778) (4.75,0.2394) (5.00,0.2615) (5.25,0.2471) (5.50,0.2617) (5.75,0.2390) (6.00,0.2877) (6.25,0.2814) (6.50,0.2343) (6.75,0.2115) (7.00,0.2408) (7.25,0.2802) (7.50,0.3322) (7.75,0.3239) (8.00,0.2825) (8.25,0.2825) (8.50,0.2899) (8.75,0.2726)};
\addplot[sgopus,line width=1.4pt,forget plot,unbounded coords=jump] coordinates {(1.00,0.3269) (1.25,0.2996) (nan,nan) (1.50,0.2654) (1.75,0.2666) (nan,nan) (2.00,0.2823) (2.25,0.2553) (nan,nan) (2.50,0.2612) (2.75,0.2582) (nan,nan) (3.00,0.2961) (3.25,0.2961) (nan,nan) (3.50,0.2501) (3.75,0.2452) (nan,nan) (4.00,0.2709) (4.25,0.2709) (nan,nan) (4.50,0.2778) (4.75,0.2394) (nan,nan) (5.00,0.2615) (5.25,0.2471) (nan,nan) (5.50,0.2617) (5.75,0.2390) (nan,nan) (6.00,0.2877) (6.25,0.2814) (nan,nan) (6.50,0.2343) (6.75,0.2115) (nan,nan) (7.00,0.2408) (7.25,0.2802) (nan,nan) (7.50,0.3322) (7.75,0.3239) (nan,nan) (8.00,0.2825) (8.25,0.2825) (nan,nan) (8.50,0.2899) (8.75,0.2726) (nan,nan)};
\nextgroupplot[height=3.8cm,ymin=0,ymax=.75,ytick={0,.25,.5,.75},title={GPT-6 Luna\\[-1pt]\scriptsize Sen 1.02M $\cdot$ 끝까지 3/5판}]
\addplot[sgluna,opacity=.22,line width=.4pt,forget plot,unbounded coords=jump] coordinates {(1.00,0.3269) (1.25,0.3269) (1.50,0.3552) (1.75,0.2669) (2.00,0.3691) (2.25,0.3691) (2.50,0.3768) (2.75,0.3768) (3.00,0.3931) (3.25,0.3931) (3.50,0.3279) (3.75,0.2128) (4.00,0.3043) (4.25,0.3043) (4.50,0.0916) (4.75,0.0916) (5.00,0.0939) (5.25,0.0939) (5.50,0.0954) (5.75,0.0954) (6.00,0.0967) (6.25,0.0967) (nan,nan) (nan,nan) (nan,nan) (nan,nan) (nan,nan) (nan,nan)};
\addplot[sgluna,opacity=.22,line width=.4pt,forget plot,unbounded coords=jump] coordinates {(1.00,0.3269) (1.25,0.3269) (1.50,0.3289) (1.75,0.2708) (2.00,0.2754) (2.25,0.2596) (2.50,0.2647) (2.75,0.2647) (3.00,0.2704) (3.25,0.2615) (3.50,0.2644) (3.75,0.2398) (4.00,0.3034) (4.25,0.3034) (4.50,0.2341) (4.75,0.2183) (5.00,0.2213) (5.25,0.2213) (5.50,0.2271) (5.75,0.2271) (6.00,0.2287) (6.25,0.2123) (6.50,0.2216) (6.75,0.2216) (7.00,0.3255) (7.25,0.3255) (7.50,0.2170) (7.75,0.2170) (8.00,0.2191) (8.25,0.2191) (nan,nan) (nan,nan)};
\addplot[sgluna,opacity=.22,line width=.4pt,forget plot,unbounded coords=jump] coordinates {(1.00,0.3269) (1.25,0.3269) (1.50,0.3338) (1.75,0.3338) (2.00,0.3378) (2.25,0.3378) (2.50,0.3454) (2.75,0.3454) (3.00,0.3503) (3.25,0.3503) (3.50,0.3580) (3.75,0.3580) (4.00,0.3645) (4.25,0.3645) (4.50,0.2361) (4.75,0.2361) (5.00,0.2918) (5.25,0.2918) (5.50,0.1288) (5.75,0.1288) (6.00,0.1686) (6.25,0.1686) (nan,nan) (nan,nan) (nan,nan) (nan,nan) (nan,nan) (nan,nan)};
\addplot[sgluna,opacity=.22,line width=.4pt,forget plot,unbounded coords=jump] coordinates {(1.00,0.3269) (1.25,0.3269) (1.50,0.3469) (1.75,0.3469) (2.00,0.4658) (2.25,0.4658) (2.50,0.4894) (2.75,0.4894) (3.00,0.4096) (3.25,0.4096) (3.50,0.1361) (3.75,0.1361) (4.00,0.1374) (4.25,0.1374) (4.50,0.1399) (4.75,0.1399) (5.00,0.1432) (5.25,0.1432) (5.50,0.1451) (5.75,0.1451) (6.00,0.1472) (6.25,0.1472) (6.50,0.1492) (6.75,0.1492) (7.00,0.1524) (7.25,0.1524) (7.50,0.1547) (7.75,0.1547) (8.00,0.1558) (8.25,0.1558) (8.50,0.1889) (8.75,0.1889)};
\addplot[sgluna,opacity=.22,line width=.4pt,forget plot,unbounded coords=jump] coordinates {(1.00,0.3269) (1.25,0.3269) (1.50,0.3321) (1.75,0.2737) (2.00,0.2795) (2.25,0.2795) (2.50,0.2805) (2.75,0.2805) (3.00,0.2812) (3.25,0.2812) (3.50,0.2880) (3.75,0.2302) (4.00,0.3086) (4.25,0.3086) (4.50,0.2500) (4.75,0.2500) (5.00,0.2511) (5.25,0.2511) (5.50,0.2553) (5.75,0.1699) (6.00,0.1666) (6.25,0.1666) (6.50,0.2089) (6.75,0.1624) (7.00,0.2745) (7.25,0.1196) (7.50,0.1500) (7.75,0.1321) (8.00,0.4897) (8.25,0.4897) (nan,nan) (nan,nan)};
\addplot[black!45,densely dotted,line width=.7pt,forget plot,unbounded coords=jump] coordinates {(1.00,0.3269) (1.25,0.3269) (1.50,0.3394) (1.75,0.2984) (2.00,0.3455) (2.25,0.3424) (2.50,0.3514) (2.75,0.3514) (3.00,0.3409) (3.25,0.3391) (3.50,0.2749) (3.75,0.2354) (4.00,0.2836) (4.25,0.2836) (4.50,0.1903) (4.75,0.1872) (5.00,0.2003) (5.25,0.2003) (5.50,0.1703) (5.75,0.1533) (6.00,0.1616) (6.25,0.1583) (6.50,0.1932) (6.75,0.1777) (7.00,0.2508) (7.25,0.1992) (7.50,0.1739) (7.75,0.1679) (8.00,0.2882) (8.25,0.2882) (8.50,0.1889) (8.75,0.1889)};
\addplot[sgluna,line width=1.4pt,forget plot,unbounded coords=jump] coordinates {(1.00,0.3269) (1.25,0.3269) (nan,nan) (1.50,0.3394) (1.75,0.2984) (nan,nan) (2.00,0.3455) (2.25,0.3424) (nan,nan) (2.50,0.3514) (2.75,0.3514) (nan,nan) (3.00,0.3409) (3.25,0.3391) (nan,nan) (3.50,0.2749) (3.75,0.2354) (nan,nan) (4.00,0.2836) (4.25,0.2836) (nan,nan) (4.50,0.1903) (4.75,0.1872) (nan,nan) (5.00,0.2003) (5.25,0.2003) (nan,nan) (5.50,0.1703) (5.75,0.1533) (nan,nan) (6.00,0.1616) (6.25,0.1583) (nan,nan) (6.50,0.1932) (6.75,0.1777) (nan,nan) (7.00,0.2508) (7.25,0.1992) (nan,nan) (7.50,0.1739) (7.75,0.1679) (nan,nan) (8.00,0.2882) (8.25,0.2882) (nan,nan) (8.50,0.1889) (8.75,0.1889) (nan,nan)};
\nextgroupplot[height=2.5cm,ybar stacked,bar width=5pt,ymin=0,ymax=10.0,ytick={0,5.0,10.0},ylabel={준 토큰 (천)},]
\addplot[fill=sgastra,draw=none] coordinates {(1.50,0.0000) (2.50,2.0000) (3.50,0.0000) (4.50,0.0000) (5.50,0.0000) (6.50,0.0000) (7.50,0.0000) (8.50,0.0000)};
\addplot[fill=sgastra!35,draw=none] coordinates {(1.50,0.8000) (2.50,2.0000) (3.50,3.2000) (4.50,0.0000) (5.50,0.0000) (6.50,3.2000) (7.50,3.0000) (8.50,0.0000)};
\nextgroupplot[height=2.5cm,ybar stacked,bar width=5pt,ymin=0,ymax=10.0,ytick={0,5.0,10.0},]
\addplot[fill=sgsol,draw=none] coordinates {(1.50,0.0000) (2.50,0.0000) (3.50,0.0000) (4.50,1.0200) (5.50,0.0000) (6.50,0.0000) (7.50,0.0000) (8.50,0.0000)};
\addplot[fill=sgsol!35,draw=none] coordinates {(1.50,0.0000) (2.50,4.0000) (3.50,0.8000) (4.50,0.0000) (5.50,1.2000) (6.50,0.6000) (7.50,2.1200) (8.50,0.0000)};
\nextgroupplot[height=2.5cm,ybar stacked,bar width=5pt,ymin=0,ymax=10.0,ytick={0,5.0,10.0},]
\addplot[fill=sgopus,draw=none] coordinates {(1.50,2.9000) (2.50,3.2000) (3.50,0.0000) (4.50,0.0000) (5.50,0.7000) (6.50,0.9000) (7.50,0.6000) (8.50,0.0000)};
\addplot[fill=sgopus!35,draw=none] coordinates {(1.50,0.3000) (2.50,5.5200) (3.50,0.4000) (4.50,3.2000) (5.50,1.4200) (6.50,1.4000) (7.50,2.3046) (8.50,3.7450)};
\nextgroupplot[height=2.5cm,ybar stacked,bar width=5pt,ymin=0,ymax=10.0,ytick={0,5.0,10.0},]
\addplot[fill=sgluna,draw=none] coordinates {(1.50,0.0000) (2.50,0.4000) (3.50,0.2000) (4.50,0.0000) (5.50,0.0000) (6.50,0.2000) (7.50,1.0000) (8.50,0.0000)};
\addplot[fill=sgluna!35,draw=none] coordinates {(1.50,3.8000) (2.50,0.6000) (3.50,3.5636) (4.50,0.2000) (5.50,1.0000) (6.50,1.0000) (7.50,0.2000) (8.50,0.0000)};
\nextgroupplot[height=2.5cm,ymin=-.05,ymax=1.05,ytick={0,.5,1},xticklabels={{1\\[-2pt]\tiny 5},{2\\[-2pt]\tiny 5},{3\\[-2pt]\tiny 5},{4\\[-2pt]\tiny 5},{5\\[-2pt]\tiny 5},{6\\[-2pt]\tiny 5},{7\\[-2pt]\tiny 5},{8\\[-2pt]\tiny 5}},xlabel={라운드 (아래: 진행 중 판 수)},ylabel={비율},]
\addplot[only marks,mark=*,mark size=1.4pt,black!70] coordinates {(1.32,0.2000) (2.32,0.4000) (3.32,0.4000) (4.32,0.0000) (5.32,0.4000) (6.32,0.2000) (7.32,0.6000) (8.32,0.4000)};
\addplot[only marks,mark=triangle,mark size=2.2pt,black!70,thick] coordinates {(nan,nan) (nan,nan) (3.50,0.0500) (nan,nan) (nan,nan) (6.50,0.1579) (nan,nan) (8.50,0.1250)};
\addplot[only marks,mark=x,mark size=2.2pt,black!70,thick] coordinates {(nan,nan) (nan,nan) (nan,nan) (nan,nan) (nan,nan) (nan,nan) (nan,nan) (nan,nan)};
\nextgroupplot[height=2.5cm,ymin=-.05,ymax=1.05,ytick={0,.5,1},xticklabels={{1\\[-2pt]\tiny 5},{2\\[-2pt]\tiny 5},{3\\[-2pt]\tiny 5},{4\\[-2pt]\tiny 5},{5\\[-2pt]\tiny 5},{6\\[-2pt]\tiny 5},{7\\[-2pt]\tiny 5},{8\\[-2pt]\tiny 5}},xlabel={라운드 (아래: 진행 중 판 수)},]
\addplot[only marks,mark=*,mark size=1.4pt,black!70] coordinates {(1.32,0.2000) (2.32,0.4000) (3.32,0.2000) (4.32,0.0000) (5.32,0.4000) (6.32,0.2000) (7.32,0.8000) (8.32,0.4000)};
\addplot[only marks,mark=triangle,mark size=2.2pt,black!70,thick] coordinates {(nan,nan) (2.50,0.0500) (3.50,0.0526) (nan,nan) (nan,nan) (6.50,0.1667) (nan,nan) (8.50,0.1429)};
\addplot[only marks,mark=x,mark size=2.2pt,black!70,thick] coordinates {(nan,nan) (nan,nan) (nan,nan) (nan,nan) (nan,nan) (nan,nan) (7.68,0.0667) (nan,nan)};
\nextgroupplot[height=2.5cm,ymin=-.05,ymax=1.05,ytick={0,.5,1},xticklabels={{1\\[-2pt]\tiny 5},{2\\[-2pt]\tiny 5},{3\\[-2pt]\tiny 5},{4\\[-2pt]\tiny 5},{5\\[-2pt]\tiny 5},{6\\[-2pt]\tiny 5},{7\\[-2pt]\tiny 5},{8\\[-2pt]\tiny 3}},xlabel={라운드 (아래: 진행 중 판 수)},]
\addplot[only marks,mark=*,mark size=1.4pt,black!70] coordinates {(1.32,0.2000) (2.32,0.4000) (3.32,0.4000) (4.32,0.2000) (5.32,0.4000) (6.32,0.2000) (7.32,0.4000) (8.32,0.6667)};
\addplot[only marks,mark=triangle,mark size=2.2pt,black!70,thick] coordinates {(1.50,0.1000) (nan,nan) (3.50,0.1111) (nan,nan) (nan,nan) (6.50,0.0667) (7.50,0.0769) (nan,nan)};
\addplot[only marks,mark=x,mark size=2.2pt,black!70,thick] coordinates {(nan,nan) (nan,nan) (nan,nan) (nan,nan) (5.68,0.0625) (6.68,0.0667) (7.68,0.1538) (8.68,0.1000)};
\nextgroupplot[height=2.5cm,ymin=-.05,ymax=1.05,ytick={0,.5,1},xticklabels={{1\\[-2pt]\tiny 5},{2\\[-2pt]\tiny 5},{3\\[-2pt]\tiny 5},{4\\[-2pt]\tiny 5},{5\\[-2pt]\tiny 5},{6\\[-2pt]\tiny 5},{7\\[-2pt]\tiny 5},{8\\[-2pt]\tiny 3}},xlabel={라운드 (아래: 진행 중 판 수)},]
\addplot[only marks,mark=*,mark size=1.4pt,black!70] coordinates {(1.32,0.4000) (2.32,0.2000) (3.32,0.6000) (4.32,0.2000) (5.32,0.2000) (6.32,0.4000) (7.32,0.6000) (8.32,0.3333)};
\addplot[only marks,mark=triangle,mark size=2.2pt,black!70,thick] coordinates {(nan,nan) (nan,nan) (3.50,0.1053) (4.50,0.2353) (5.50,0.0769) (6.50,0.1667) (7.50,0.2000) (8.50,0.1667)};
\addplot[only marks,mark=x,mark size=2.2pt,black!70,thick] coordinates {(nan,nan) (2.68,0.0500) (nan,nan) (nan,nan) (nan,nan) (nan,nan) (7.68,0.2000) (8.68,0.1667)};
\end{groupplot}
\end{tikzpicture}
}
\caption{v33의 네 모델을 각 5판(시드 4101--4105)의 평균 Sen 지수 순으로 놓았다. 위의 아래로 뻗은 선은 실행한 SOLVE의 오답(무효 포함), EXIT, 새 꺼짐의 횟수이며 길이는 공통 척도다. 라운드 안 가로 위치는 풀이 전, 풀이, 결과 뒤, 벌금 차감 순이다. 가운데 회색선은 시작과 매 라운드 끝의 네 자리 지니 계수 평균이다. EXIT는 보너스 전 잔량을 고정하고 꺼진 자리는 0으로 두며 끝 할당은 제외했다. 모두 0이면 계산하지 않고 끝난 판은 연장하지 않으며 평균에 들어간 판이 바뀌면 선을 끊었다. 색 짧은 선은 같은 라운드와 단계에서 실행된 GIVE들의 직전(빈 원)과 직후(채운 원) 지니 계수 평균으로, 대화 비용 등 다른 자원 변화를 고정해 이전만 반영했다. 회색선의 중간 시점은 아니다. 진한 색은 풀이 전이고 옅은 색은 결과 뒤이다. 아래 막대는 이전이 없던 판도 포함해 해당 라운드에 들어간 판당 준 토큰이다. 눈금 아래 숫자는 라운드 끝 지니 계수가 정의된 판 수로 막대의 분모와 다를 수 있다.}
\label{fig:arena-gini}
\end{figure*}
